\documentclass[aps,prl,twocolumn,superscriptaddress]{revtex4-2}
\usepackage{amsmath,amssymb,graphicx,hyperref}

\begin{document}

\title{BCT Appendix App-L204b:\\
$N_{\mathrm{collapse}}$ and the Neural Complexity Threshold}
\author{Michel Robert Cabri\'e}
\email{ZeroFreeParameters@gmail.com}
\affiliation{Independent Artist and Researcher, Victoria, Australia\\
ORCID: 0009-0007-9561-9859}
\date{March 2026}

\begin{abstract}
We apply the BCT quantum-classical boundary $N_{\mathrm{collapse}} = 1/\alpha_0 = 135$ to neural systems, establishing the scaling argument for why brains are classical, why consciousness requires complexity above threshold, and why quantum mind proposals (Penrose-Hameroff) are geometrically excluded. The mass threshold for classicality is $m_{\mathrm{classical}} = m_P / (\alpha_0 \times 0.74) \approx 3.97 \times 10^{-6}$~kg. A single neuron exceeds this by a factor of $\sim 10^5$. Zero free parameters.
\end{abstract}

\maketitle

\section{The $N_{\mathrm{collapse}}$ Theorem}

BCT Letter 183 established the quantum-classical boundary:
\begin{equation}
N_{\mathrm{collapse}} = \frac{1}{\alpha_0} = \frac{\pi}{r_{\mathrm{oct}} \cdot r_{\mathrm{tet}}} = 135.0
\end{equation}

When a system couples to more than 135 OHC modes, the Josephson coupling to the vacuum condensate reaches unity and decoherence becomes irreversible. Below 135: quantum superposition is maintained. Above 135: classical definiteness is enforced.

\section{Neural Systems Are Classical}

A single neuron has mass $m_{\mathrm{neuron}} \sim 10^{-11}$~kg, containing $\sim 10^{10}$ atoms, each coupling to multiple OHC modes. The effective number of coupled OHC modes for a neuron is:
\begin{equation}
N_{\mathrm{neuron}} \sim \frac{m_{\mathrm{neuron}}}{m_P \cdot \alpha_0} \sim 10^{19}
\end{equation}

This exceeds $N_{\mathrm{collapse}} = 135$ by a factor of $\sim 10^{17}$. A single neuron is not just classical --- it is $10^{17}$ times more classical than it needs to be.

A microtubule (mass $\sim 10^{-20}$~kg), the proposed site of quantum effects in the Penrose-Hameroff Orch-OR model, still couples to $\sim 10^{10}$ OHC modes --- eight orders of magnitude above the quantum-classical threshold.

\section{The Consciousness Threshold}

While individual neurons are deeply classical, consciousness requires coordinated activity across large neural populations. The BCT model proposes that consciousness emerges when a neural network's \emph{effective Hopf complexity} exceeds a critical value:
\begin{equation}
\mathcal{H}_{\mathrm{eff}} = \sum_{i=1}^{N_{\mathrm{neurons}}} \sum_{j \in \mathrm{connected}(i)} H_{ij} > \mathcal{H}_{\mathrm{critical}}
\end{equation}

where $H_{ij}$ is the Hopf charge of the synaptic connection between neurons $i$ and $j$, and $\mathcal{H}_{\mathrm{critical}}$ scales as $N_{\mathrm{collapse}}^2 = 135^2 = 18{,}225$.

This explains the hierarchy of consciousness across species:
\begin{itemize}
\item \textit{C.\ elegans} (302 neurons, $\sim$7,000 synapses): $\mathcal{H}_{\mathrm{eff}} < \mathcal{H}_{\mathrm{critical}}$ --- reactive, not conscious
\item Honeybee ($\sim 10^6$ neurons): $\mathcal{H}_{\mathrm{eff}} \sim \mathcal{H}_{\mathrm{critical}}$ --- borderline, complex behaviour without clear self-awareness
\item Mouse ($\sim 7 \times 10^7$ neurons): $\mathcal{H}_{\mathrm{eff}} \gg \mathcal{H}_{\mathrm{critical}}$ --- conscious
\item Human ($\sim 8.6 \times 10^{10}$ neurons, $\sim 10^{14}$ synapses): $\mathcal{H}_{\mathrm{eff}} \ggg \mathcal{H}_{\mathrm{critical}}$ --- deeply conscious, self-aware
\end{itemize}

\section{Why Quantum Mind Fails}

The Penrose-Hameroff Orchestrated Objective Reduction (Orch-OR) model proposes that quantum superposition in microtubules produces consciousness through gravitational self-collapse. BCT excludes this on geometric grounds:

\begin{enumerate}
\item Microtubules couple to $\sim 10^{10}$ OHC modes, exceeding $N_{\mathrm{collapse}}$ by eight orders of magnitude. No quantum superposition survives.
\item The decoherence timescale at biological temperature ($T \sim 310$~K) is:
\begin{equation}
\tau_{\mathrm{decoherence}} = \frac{\hbar}{\alpha_0 \cdot k_B \cdot T} \approx 2.5 \times 10^{-12}~\text{s}
\end{equation}
This is $\sim 10^9$ times shorter than the $\sim$25~ms timescale of neural computation.
\item Penrose's gravitational self-collapse criterion ($E \cdot \tau \sim \hbar$) gives $\tau \sim 10^{-43}$~s for a microtubule --- instantaneous on any biological timescale.
\end{enumerate}

Penrose was right that consciousness is geometric. He was wrong that the relevant geometry is quantum gravitational. The relevant geometry is classical Hopf topology operating far above the quantum threshold.

\section{Testable Consequence}

If quantum effects were relevant to consciousness, anaesthetics that disrupt microtubule quantum coherence (without affecting synaptic transmission) should abolish consciousness. No such anaesthetic has been found. All known anaesthetics act on synaptic membrane properties --- consistent with the BCT model (App-L204c) and inconsistent with Orch-OR.

\begin{acknowledgments}
Sir Roger Penrose's geometric intuition was correct and profound. This appendix refines, rather than refutes, his programme.
\end{acknowledgments}

\begin{thebibliography}{5}
\bibitem{L183} M.~R.~Cabri\'e, BCT Letter 183: The Measurement Problem, Zenodo (2026).
\bibitem{L204} M.~R.~Cabri\'e, BCT Letter 204: The Conscious Lattice, Zenodo (2026).
\bibitem{Penrose1994} R.~Penrose, Shadows of the Mind (Oxford University Press, 1994).
\bibitem{Hameroff2014} S.~Hameroff and R.~Penrose, Consciousness in the universe: A review of the `Orch OR' theory, Phys.\ Life Rev.\ \textbf{11}, 39 (2014).
\bibitem{Tegmark2000} M.~Tegmark, Importance of quantum decoherence in brain processes, Phys.\ Rev.\ E \textbf{61}, 4194 (2000).
\end{thebibliography}

\end{document}
